\documentclass[UTF8,fontset=none,11pt,a4paper]{ctexart}
\usepackage{ipara-handbook}
\handbooksetup{DS-11 内部排序}{408 · 数据结构 DS-11}{Progress · Invariant · Stability · Cost}
\begin{document}
\handbooktitle{内部排序}{在内存中怎样用局部进展制造全局有序}{408 · 数据结构 Topic DS-11}

\begin{corebox}{母问题：每一步究竟把什么“已排序”变成事实？}
内部排序的生成链是
\[
\text{Local Operation}\to\text{Progress Measure}\to\text{Ordering Invariant}\to\text{Termination}\to\text{Cost/Stability}.
\]
插入排序维护已排序前缀，选择排序维护已确定前缀，交换排序用 partition 缩小未决区间，归并排序把有序子段合并，堆排序借用 DS05 的根极值。算法区别不只在口诀，而在每轮不变量和数据搬运方式。
\end{corebox}
\begin{methodbox}{本册定位与 Owner 边界}
\begin{itemize}
\item \kw{Owns：}直接/折半插入、Shell、冒泡、简单选择、快速、堆、二路归并、基数排序，以及比较/移动、稳定性、原地性与最坏输入。
\item \kw{Uses：}DS05 堆的接口与 Atlas Foundation 成本语言。
\item \kw{Stops：}外部存储块 I/O 与多路归并归 DS12；选择哪种排序的 workload 规则进入 Rules。
\end{itemize}
\end{methodbox}
\tableofcontents\newpage

\section{排序契约与三种性质}
排序先固定升序/降序、是否允许重复、是否要求稳定以及是否允许额外空间。稳定性保持相等 key 的原始相对次序；原地性描述额外空间，不等于稳定；比较排序的下界适用于只通过比较获得次序的算法，基数排序改变了信息模型。

\section{不变量驱动的四条主线}
插入排序：前缀有序，取下一个元素向左移动直到放置点；最好 $O(n)$、最坏 $O(n^2)$，稳定。选择排序：前缀是全局最小的若干元素，交换次数少但通常不稳定。快速排序：partition 后枢轴落在最终位置，递归处理两侧；平均 $O(n\log n)$、最坏 $O(n^2)$。归并排序：两个有序段按双指针合并，稳定且 $O(n\log n)$，需要线性辅助空间。
\begin{examplebox}{为什么“有序区”不是一句空话}
对数组 $[4,1,3,2]$ 做插入排序。处理完 $4,1$ 后前缀成为 $[1,4]$；处理 $3$ 时只需在该前缀内找插入点。若每轮都重新排序整个数组，就失去该不变量带来的结构化进展。
\end{examplebox}

\section{插入族：定位成本与移动成本必须分开}
直接插入从有序前缀末尾向左比较并移动，直到找到插入位置。折半插入只用二分缩短“找位置”的比较次数，但连续数组仍要移动插入点后的元素，所以总移动量和最坏时间仍为 $O(n^2)$；它不是把插入排序整体变成 $O(n\log n)$。

Shell 排序选择递减增量 $d_1>d_2>\cdots>1$，每一趟对所有相距 $d$ 的子序列做插入排序。大增量先让元素跨远距离移动，最后 $d=1$ 时数组已经较接近有序。其正确性最终由最后一趟普通插入保证；复杂度依赖增量序列，不能在未给序列时统一宣称精确阶。跨距移动可能改变相等 key 的相对次序，因此 Shell 通常不稳定。

\section{交换与选择族：每趟提交的事实不同}
冒泡排序比较相邻逆序对并交换；一趟从左到右可把当前最大元素提交到未排序区末端。若整趟没有交换即可提前停止，因此已排序输入最好为 $O(n)$。简单选择排序每趟在未排序区找最小元素，与边界位置交换；它只进行 $O(n)$ 次交换，但一次跨越交换可能越过相等 key，通常不稳定。

快速排序用 partition 把区间分成枢轴两侧，并让至少一个分割边界成为进展。正确性依赖 partition 循环不变量；复杂度依赖分割是否均衡。已排序、逆序或大量重复 key 在某些固定枢轴实现中会持续产生极端分割，使递归深度和时间退化到 $O(n^2)$。

\section{堆、二路归并与基数：改变候选组织方式}
堆排序先调用 DS05 建立最大堆，每轮把根极值与未排序区末端交换，缩小堆边界并对新根下滤。已排序后缀不断扩大，时间稳定为 $O(n\log n)$、额外空间 $O(1)$，但根与末端的远距离交换使其不稳定。

二路归并把两个已有序段合成一个有序段：比较两段当前首元素，取较小者并推进对应指针；某段耗尽后复制另一段剩余部分。若相等时优先取左段，可保持稳定。自底向上的归并把长度 $1,2,4,\ldots$ 的段逐轮合并，无论初始次序如何都为 $O(n\log n)$，代价是 $O(n)$ 辅助数组。

基数排序不比较完整 key，而是按某一位的桶序反复稳定分配与收集。最低位优先从低位到高位，后一趟必须稳定，才能保留前面低位已建立的次序。若有 $d$ 位、基数为 $r$，成本为 $O(d(n+r))$；它依赖 key 可分解为有限位，不受比较排序下界约束。

\begin{methodbox}{一张表不能替代机制}
遇到排序过程题，先写本趟提交的不变量：插入是有序前缀，冒泡是末端极值，选择是前端极值，快速是分区边界，堆是根极值加堆边界，归并是有序段，基数是已处理位上的稳定次序。然后才数比较、移动、趟数和空间。
\end{methodbox}

\section{代码与测试证据}
完整实现位于 \codepath{code/ds11_internal_sort.hpp}，测试位于 \codepath{code/ds11_internal_sort_test.cpp}。
\lstinputlisting[language=C++,firstline=10,lastline=86,firstnumber=10,numbers=left,caption={插入、归并与快速排序核心转移}]{30_408/10_数据结构/11_内部排序/code/ds11_internal_sort.hpp}
\lstinputlisting[language=C++,firstline=8,lastline=42,firstnumber=8,numbers=left,caption={空、重复、已排序与逆序边界测试}]{30_408/10_数据结构/11_内部排序/code/ds11_internal_sort_test.cpp}
\begin{lstlisting}[language=bash]
clang++ -std=c++17 -Wall -Wextra -Werror -pedantic code/ds11_internal_sort_test.cpp -o /tmp/ds11_test
/tmp/ds11_test
\end{lstlisting}

\section{复杂度与概念边界}
\begin{center}\small
\begin{tabularx}{\linewidth}{L{2.4cm}C{1.7cm}C{1.7cm}C{1.4cm}YY}\toprule
算法 & 最好 & 平均 & 最坏 & 稳定/空间\\\midrule
直接/折半插入 & $n$ & $n^2$ & $n^2$ & 稳定、原地\\
冒泡 & $n$ & $n^2$ & $n^2$ & 稳定、原地\\
简单选择 & $n^2$ & $n^2$ & $n^2$ & 不稳定、原地\\
Shell & 依增量 & 依增量 & 依增量 & 不稳定、原地\\
归并 & $n\log n$ & $n\log n$ & $n\log n$ & 稳定、$O(n)$\\
快速 & $n\log n$ & $n\log n$ & $n^2$ & 通常不稳定\\
堆 & $n\log n$ & $n\log n$ & $n\log n$ & 不稳定、原地\\
基数 & $d(n+r)$ & $d(n+r)$ & $d(n+r)$ & 稳定分配、$O(n+r)$\\\bottomrule\end{tabularx}\end{center}

\section{做题控制协议}
先写“哪一段已排序/哪个枢轴已归位”，再数比较、移动和辅助空间。看到稳定性要求优先排除会跨越相等 key 的交换；看到最坏保证与线性空间权衡，比较归并、堆与快速的不同前提。DS05 只提供堆操作，堆排序的整体不变量由本册负责。
\begin{corebox}{压缩复原}
五问：已确定区域在哪里？一次局部操作扩大了什么？终止量是什么？相等 key 的次序是否保留？最坏输入会让哪一支退化？
\end{corebox}
\end{document}
